\documentclass[11pt]{article}
\usepackage[margin=1in]{geometry}
\usepackage{amsmath,amssymb,amsthm}
\usepackage[colorlinks=true,linkcolor=blue,urlcolor=blue,citecolor=blue]{hyperref}
\newtheorem{theorem}{Theorem}[section]
\newtheorem{lemma}[theorem]{Lemma}
\newtheorem{corollary}[theorem]{Corollary}
\newcommand{\R}{\mathbb R}
\newcommand{\norm}[1]{\left\lVert#1\right\rVert}
\newcommand{\tr}{\operatorname{tr}}
\newcommand{\rank}{\operatorname{rank}}
\newcommand{\diag}{\operatorname{diag}}
\newcommand{\spec}{\operatorname{spec}}
\title{Every determinant-admissible relative spectrum\protect\\
can arise from one edge change}
\author{Research working manuscript}
\date{November 2026}
\begin{document}
\maketitle
\begin{abstract}
Fix $0<\alpha<1$ and a rank-one positive update with fixed leverage
$u^TA^{-1}u=\tau$. We characterize the closure of the relative spectra
of $(A+uu^T)^\alpha$ with respect to $A^\alpha$: positivity and the
determinant identity are the only restrictions. A separated-scale
construction realizes every admissible spectrum in the limit.
Consequently the exact supremal number of relative correction
eigenvalues exceeding $\varepsilon$ is the largest nonnegative integer
strictly less than $\alpha\log(1+\tau)/\log(1+\varepsilon)$.
The same supremum already occurs for a single edge-conductance increase
on connected weighted paths with ordinary unit vertex masses. In
particular a bounded effective-resistance budget permits relative
compression rank of order $1/\varepsilon$, despite the logarithmic-squared
absolute compression bounds available for fractional updates. The upper
bound is an elementary determinant consequence; the substantive assertion
is full sharpness and its realization on unit-mass paths. Historical
priority has not been certified.
\end{abstract}

\section{The relative spectral question}
All matrices are real and finite dimensional. For $A\succ0$,
$B=A+uu^T$, and $0<\alpha<1$, define
\[
 Q_\alpha(A,u)=A^{-\alpha/2}B^\alpha A^{-\alpha/2},\qquad
 C_\alpha(A,u)=Q_\alpha(A,u)-I.
\]
The leverage is $\tau=u^TA^{-1}u$. Powers are principal powers.
The eigenvalues of $Q_\alpha$ are the generalized eigenvalues of
$(B^\alpha,A^\alpha)$. This is an energy-relative comparison; it is
different from dividing $\norm{B^\alpha-A^\alpha}$ by its own norm.

For connected graph Laplacians, all inverse powers in this paper mean
the Moore--Penrose functional calculus: zero on the common kernel and
the stated negative power on its orthogonal complement. The relative
matrix $C$ is zero on that kernel. Every determinant on the common
support is explicitly understood as such.

\begin{theorem}[Complete closure of rank-one relative spectra]
\label{thm:closure}
Fix an integer $n\ge1$, $0<\alpha<1$, and $\tau\ge0$.
The closure of the ordered eigenvalue lists of
$Q_\alpha(A,u)$ over all $A\in\R^{n\times n}$ positive definite and
all $u\in\R^n$ with $u^TA^{-1}u=\tau$ is exactly
\[
 \mathcal S_{n,\alpha,\tau}
 =\left\{(q_1,\ldots,q_n):
       q_1\ge\cdots\ge q_n\ge1,\quad
       \prod_{i=1}^nq_i=(1+\tau)^\alpha\right\}.
\]
For $\tau=0$ this set is the singleton $(1,\ldots,1)$.
The closure is essential: the theorem asserts approximation of each
list, without asserting that every list is attained at finite condition
number.
\end{theorem}

\begin{corollary}[Exact worst-case strict count]\label{cor:count}
For $\tau>0$ and $\varepsilon>0$, put
\[
 R=\frac{\alpha\log(1+\tau)}{\log(1+\varepsilon)},\qquad
 M=\max\{0,\lceil R\rceil-1\}.
\]
Then
\[
 \sup_{A\succ0,\ u^TA^{-1}u=\tau}
 \#\{j:\lambda_j(C_\alpha(A,u))>\varepsilon\}=M,
\]
where the supremum ranges over all finite dimensions.
In a fixed dimension $n$, the exact supremum is $\min(n,M)$.
For $\tau=0$ both suprema are zero.
\end{corollary}

For fixed $\tau>0$, the worst-case count is
\[
 M=\frac{\alpha\log(1+\tau)}{\varepsilon}+O_{\alpha,\tau}(1)
 \quad(\varepsilon\downarrow0).
\]
Thus no uniform logarithmic-in-accuracy rank bound is possible in this
relative metric, even for a single positive rank-one update of fixed
leverage. This does not contradict absolute-error compression estimates.

\section{Proof of the spectral closure}
The following integral supplies the whole construction.
Write $c_\alpha=\sin(\pi\alpha)/\pi>0$.

\begin{lemma}[Separated-scale rank-one limit]\label{lem:scales}
Let $r_1,\ldots,r_n\ge1$, and set
\[
 z_i=(r_i-1)\prod_{j=i+1}^n r_j,\qquad
 a_i=\delta^{n-i},\quad
 A_\delta=\diag(a_1,\ldots,a_n),\quad
 (u_\delta)_i=\sqrt{a_i z_i},\qquad 0<\delta<1.
\]
Then, for every $\delta$,
\[
 u_\delta^TA_\delta^{-1}u_\delta
    =\sum_i z_i=\prod_i r_i-1.
\]
Moreover, in operator norm as $\delta\downarrow0$,
\[
 C_\alpha(A_\delta,u_\delta)
       \longrightarrow\diag(r_1^\alpha-1,\ldots,r_n^\alpha-1).
\]
\end{lemma}
\begin{proof}
The telescoping identity for $\sum_i z_i$ is immediate from
$z_i=\prod_{j=i}^n r_j-\prod_{j=i+1}^n r_j$.
The resolvent formula for fractional powers and the rank-one inverse
identity give
\begin{equation}\label{eq:resolvent}
 B^\alpha-A^\alpha
 =c_\alpha\int_0^\infty
 \frac{t^\alpha(A+tI)^{-1}uu^T(A+tI)^{-1}}
 {1+u^T(A+tI)^{-1}u}\,dt.
\end{equation}
The integral converges in norm; its integrand is positive semidefinite.
For its $i$th relative diagonal entry, substitute $t=a_i s$ to obtain
\begin{equation}\label{eq:diag}
 C_{ii}=c_\alpha z_i\int_0^\infty
 \frac{s^\alpha}{(1+s)^2}
 \left(1+\sum_j\frac{a_jz_j}{a_j+a_i s}\right)^{-1}ds.
\end{equation}
For every $s>0$, $a_j/a_i\to0$ if $j<i$, and $a_j/a_i\to\infty$
if $j>i$. The denominator in parentheses therefore tends to
\[
 H_i+\frac{z_i}{1+s},\qquad
 H_i=1+\sum_{j>i}z_j=\prod_{j>i}r_j.
\]
The integrable function $s^\alpha/(1+s)^2$ dominates the integrand
after removing the fixed multiplier $c_\alpha z_i$. Dominated
convergence gives
\[
 C_{ii}\longrightarrow
 c_\alpha\frac{z_i}{H_i}
    \int_0^\infty\frac{s^\alpha}{(1+s)(1+s+z_i/H_i)}\,ds
 = (1+z_i/H_i)^\alpha-1=r_i^\alpha-1.
\]
The scalar equality follows by applying the same fractional resolvent
formula to the scalars $1+z_i/H_i$ and $1$. It also holds when $z_i=0$.

For $i<j$, omit the positive denominator in \eqref{eq:resolvent}.
The resulting scalar divided-difference integral yields
\begin{align*}
 |C_{ij}|
 &\le\sqrt{z_i z_j}(a_i a_j)^{(1-\alpha)/2}
                  \frac{a_j^\alpha-a_i^\alpha}{a_j-a_i}\\
 &=\sqrt{z_i z_j}\,x^{(1-\alpha)/2}
                    \frac{1-x^\alpha}{1-x},\qquad x=a_i/a_j.
\end{align*}
This tends to zero because $0<\alpha<1$ and $x\to0$.
Symmetry treats $i>j$. Entrywise convergence in fixed dimension implies
operator-norm convergence.
\end{proof}

\begin{proof}[Proof of Theorem~\ref{thm:closure}]
Formula \eqref{eq:resolvent} proves $B^\alpha-A^\alpha\succeq0$,
so every eigenvalue of $Q_\alpha$ is at least one. The determinant
lemma gives
\[
 \det Q_\alpha
    =\left(\frac{\det(A+uu^T)}{\det A}\right)^\alpha
    =(1+\tau)^\alpha.
\]
These constraints are closed. Conversely, given an admissible list,
take $r_i=q_i^{1/\alpha}$ and apply Lemma~\ref{lem:scales}.
Its leverage is exactly $\tau$ at every $\delta$, and its relative
eigenvalues converge to the desired list, after ordering.
For $\tau=0$, positive definiteness forces $u=0$.
\end{proof}

\begin{proof}[Proof of Corollary~\ref{cor:count}]
If $m$ correction eigenvalues exceed $\varepsilon$, then
\[
 (1+\varepsilon)^m<\prod_j(1+\lambda_j(C))
                 =(1+\tau)^\alpha,
\]
so $m<R$. The largest allowable integer is $M$; in dimension $n$ it
is also at most $n$.
For the reverse inequality, let $m=\min(n,M)>0$. Choose a target
list with $m$ entries equal to $(1+\tau)^{\alpha/m}>1+\varepsilon$
and the remaining $n-m$ entries equal to one. Theorem~\ref{thm:closure}
approximates it closely enough to retain those $m$ strict inequalities.
For the unrestricted supremum take dimension $M$ when $M>0$.
The case $M=0$ is immediate. At integral $R$, the strict inequality
excludes $m=R$, which explains $\lceil R\rceil-1$ rather than
$\lfloor R\rfloor$.
\end{proof}

\section{Transfer to one edge on an ordinary weighted path}
For a connected weighted graph with Laplacian $L$, write
\[
 \mathcal C_\alpha(L,L')
  =(L^\dagger)^{\alpha/2}(L'^\alpha-L^\alpha)(L^\dagger)^{\alpha/2}.
\]
All graph conductances below are strictly positive on path edges.
The update increases one existing edge of conductance $w_0$ by
$w_1=\tau w_0$. On a path this is equivalent to a fixed leverage
$w_1 b^TL^\dagger b=\tau$, because the effective resistance across a
bridge is $1/w_0$.

\begin{theorem}[Unit-mass weighted-path sharpness]\label{thm:path}
Fix $0<\alpha<1$ and $\tau>0$. For every finite list
$q_i\ge1$ with $\prod_iq_i=(1+\tau)^\alpha$, and every tolerance
$\eta>0$, there is a finite connected weighted path with ordinary
unit vertex masses and a single edge-conductance increase by the
factor $1+\tau$ such that the eigenvalues of
$\mathcal C_\alpha(L,L')$ are, after ordering, within $\eta$ of
\[
 (q_1-1,\ldots,q_n-1,0,\ldots,0).
\]
The number of trailing zeros is allowed to grow with the construction;
the assertion compares finite lists padded to the constructed dimension.
Consequently the exact unrestricted strict-count supremum on this class
of ordinary weighted paths is $M$ from Corollary~\ref{cor:count}.
\end{theorem}

The path theorem uses variable positive conductances and unrestricted
path length. It does not assert the same law for unweighted paths, for
bounded conductance ratios, or in a fixed path dimension.

\begin{lemma}[Mass path realization]\label{lem:masspath}
The relative eigenvalue list of every SPD rank-one witness with
$\tau>0$, padded with one zero correction eigenvalue, is a limit of
relative eigenvalue lists of one-edge updates on connected weighted
paths with positive vertex masses. Its leverage is preserved exactly
throughout. Matrix norm convergence is understood after orthogonal
identification of the eigenbases.
\end{lemma}
\begin{proof}
It suffices first to treat a cyclic pair $(A,u)$ with $u\ne0$.
Orthogonal Lanczos, followed by alternating basis signs, transforms it
into $(T,\sqrt\kappa e_1)$, where $\kappa=\norm u^2$ and $T$ is SPD
tridiagonal with strictly negative off-diagonal entries.
Let $d=T^{-1}e_1$. All entries of $d$ are positive: write
$T=sI-N$ with $N\ge0$ irreducible and $\norm N<s$ for sufficiently
large $s$, and expand its inverse in a geometric series. Put
$D=\diag(d)$ and $K_0=DTD$. Then
\[
 K_0\mathbf1=d_1e_1.
\]
Thus $K_0$ is the conductance matrix of a path grounded at vertex one,
with ground-edge conductance $d_1$. Add the ground vertex with mass
$R^2$, and give internal vertices masses $d_i^2$. Its mass-normalized
Laplacian and updated Laplacian are
\[
 H_R=\begin{pmatrix}d_1/R^2&-e_1^T/R\\-e_1/R&T\end{pmatrix},
 \qquad H_R'=H_R+\kappa v_Rv_R^T,
 \qquad v_R=(-d_1/R,e_1)^T.
\]
The physical ground-edge increase is $\kappa d_1^2$, and hence its
leverage is exactly $\kappa d_1=\kappa e_1^TT^{-1}e_1=\tau$.
Both endpoints have the same kernel spanned by $(R,d)^T$.
As $R\to\infty$,
\[
 H_R\to0\oplus T,\qquad
 H_R'\to0\oplus(T+\kappa e_1e_1^T).
\]
There is exactly one zero eigenvalue and all other eigenvalues stay
bounded below by a positive constant for sufficiently large $R$.
Consequently both the positive fractional calculus and the
pseudoinverse negative fractional calculus converge in norm. Their
relative correction converges to $0\oplus C_\alpha(T,\sqrt\kappa e_1)$.

For a noncyclic witness, choose SPD matrices $A_j\to A$ with simple
spectra and cyclic vectors $\widetilde u_j\to u$. Rescale each vector
by $\sqrt{\tau/(\widetilde u_j^TA_j^{-1}\widetilde u_j)}$; these
positive factors tend to one. The rescaled cyclic witnesses have
exactly the original leverage and their relative matrices converge to
the original relative matrix by continuous functional calculus.
Apply the cyclic construction and take a diagonal sequence of finite
path approximations. Every path has the original witness dimension
plus one.
\end{proof}

\begin{lemma}[Removing the vertex masses]\label{lem:blowup}
Let $K,K'$ be connected conductance Laplacians of a fixed weighted
path differing at one edge, and let $M=\diag(M_i)\succ0$ be fixed
vertex masses. The relative correction for
$H=M^{-1/2}KM^{-1/2}$ and $H'=M^{-1/2}K'M^{-1/2}$ is a norm limit,
with extra zero eigenvalues, of relative corrections for ordinary
unit-mass weighted paths. The edge increase ratio is unchanged.
\end{lemma}
\begin{proof}
Choose integers $m_i(s)\ge1$ with $m_i(s)/s\to M_i$.
For each fixed integer $s$, replace vertex $i$ by a consecutive block
of $m_i(s)$ unit-mass vertices. Connect successive vertices within
each block by edges of conductance $T$, and connect consecutive blocks
by single edges of conductance $s$ times the original path conductance.
At the edited interblock edge, use $s$ times its original increment.
The resulting graph is again one connected path.

For fixed $s$, write its two Laplacians as
\[
 L_T=TS+K_{\rm ext},\qquad
 L_T'=TS+K_{\rm ext}',
\]
where $S$ is the direct sum of unit-conductance path Laplacians
inside the blocks. Let $J$ be the isometry whose $i$th column is
$m_i(s)^{-1/2}$ times the indicator of block $i$.
Then $\ker S=\operatorname{range}J$, and
\[
 J^TK_{\rm ext}J
  =\diag(m_i(s))^{-1/2}(sK)\diag(m_i(s))^{-1/2}=H_s,
\]
with the analogous identity for $H_s'$.
In the orthogonal decomposition $\operatorname{range}J\oplus
\operatorname{range}J^\perp$, the high block is $TS_1+O(1)$ with
$S_1\succ0$, the low block is $H_s$, and the off-diagonal blocks are
$O(1)$. A Schur-complement inverse therefore gives, for every fixed
$t>0$,
\[
 (L_T+tI)^{-1}\longrightarrow J(H_s+tI)^{-1}J^T,
\]
and the same holds for $L_T'$.

Put $E=L_T'-L_T$, which is fixed as $T$ varies. The resolvent
difference $D_T(t)$ obeys
\[
 \norm{D_T(t)}\le\min(2/t,\norm E/t^2).
\]
This makes $t^\alpha D_T(t)$ integrable uniformly over $T$.
The fractional difference integral and dominated convergence imply
\[
 L_T'^\alpha-L_T^\alpha
       \longrightarrow J(H_s'^\alpha-H_s^\alpha)J^T
       \quad\hbox{in operator norm}.
\]
Standard finite-dimensional block spectral separation also gives
\[
 (L_T^\dagger)^{\alpha/2}
       \longrightarrow J(H_s^\dagger)^{\alpha/2}J^T.
\]
For completeness, there are exactly $\dim\ker S$ bounded
eigenvalues, converging with their low spectral subspace to the
eigenvalues of $H_s$; all remaining eigenvalues diverge linearly in
$T$. The unique zero eigenvector is constant. The other bounded
eigenvalues converge to strictly positive numbers because the quotient
path is connected. Applying $x\mapsto x^{-\alpha/2}$ on those
positive eigenvalues, zero on the kernel, and zero in the limit on
the divergent eigenvalues proves the displayed convergence.
Multiplying the three convergent factors proves relative-correction
convergence for fixed $s$, with extra zeros on the high subspace.

Finally $H_s\to H$ and $H_s'\to H'$ because $m_i(s)/s\to M_i$.
Their ranks stay equal to the quotient size minus one, with their
positive eigenvalues bounded away from zero. The same functional
calculus continuity proves relative-correction convergence as
$s\to\infty$. These are sequential limits: first choose sufficiently
large finite $s$, then sufficiently large finite $T$.
The edited edge is a bridge in every constructed path. Both its base
conductance and increment were multiplied by $s$, so its leverage
and increment ratio are preserved exactly, independently of $T$.
\end{proof}

\begin{proof}[Proof of Theorem~\ref{thm:path}]
Combine Theorem~\ref{thm:closure}, Lemma~\ref{lem:masspath}, and
Lemma~\ref{lem:blowup}, allocating the desired tolerance among the
three finite approximation steps. Lists with entries equal to one
are covered by the noncyclic approximation in
Lemma~\ref{lem:masspath}. Every approximation uses exactly the same
leverage $\tau$. The number of extra eigenvalues is allowed to grow,
and all lie as close to zero as required.

For the strict-count upper bound restrict an ordinary connected path
to $\mathbf1^\perp$ and apply Corollary~\ref{cor:count}.
For the lower bound choose its equal-eigenvalue target strictly above
$\varepsilon$ and approximate closely enough to preserve all strict
inequalities. An additional eigenvalue exceeding $\varepsilon$ is
excluded by the already proved universal upper bound. Thus the exact
supremum is $M$.
\end{proof}

\section{Consequences and limits}
\begin{corollary}[Updated-energy version]\label{cor:updated}
Let
\[
 D_\alpha=B^{-\alpha/2}(B^\alpha-A^\alpha)B^{-\alpha/2}.
\]
If $q_j$ are the relative eigenvalues in Theorem~\ref{thm:closure},
then the eigenvalues of $D_\alpha$ are $1-q_j^{-1}$.
For $0<\varepsilon<1$, its exact unrestricted strict-count supremum
at fixed $\tau>0$ is
\[
 \max\left\{0,\left\lceil
       \frac{\alpha\log(1+\tau)}{-\log(1-\varepsilon)}
                              \right\rceil-1\right\}.
\]
This is sharp on ordinary unit-mass weighted paths with one edge
increase, as above.
\end{corollary}
\begin{proof}
The eigenvalues of $B^{-\alpha/2}A^\alpha B^{-\alpha/2}$ are the
reciprocals of those of $A^{-\alpha/2}B^\alpha A^{-\alpha/2}$.
Thus $\lambda(D)>\varepsilon$ is equivalent to
$q>1/(1-\varepsilon)$; apply the count theorem and the path transfer.
\end{proof}

For a PSD addition of any rank, $B=A+UU^T$, the elementary exact budget is
\[
 \sum_j\log(1+\lambda_j(C_\alpha))
       =\alpha\log\det(I+U^TA^{-1}U).
\]
The rank-one construction proves that even the most concentrated
possible original relative update can distribute this budget in an
arbitrary admissible manner after a fractional power.

In graph terms the best rank-$r$ energy-relative correction is the
spectral truncation of $\mathcal C_\alpha$. If its omitted eigenvalues
are at most $\varepsilon$, the reconstructed correction $F$ satisfies
\[
 0\preceq (L'^\alpha-L^\alpha)-F\preceq\varepsilon L^\alpha
 \quad\hbox{on }\mathbf1^\perp.
\]
The count theorem gives the exact worst-case rank required by this
criterion. This is a structural existence and obstruction theorem;
no nearly-linear algorithm for this relative truncation is proved here.
The constructed path sizes and conductance ratios can be arbitrarily
large as the approximation tolerance shrinks.

\section{Relation to existing matrix-update work}
Low-rank rational matrix-function updates and their graph applications
are established~\cite{BCKS}. Square-root update decay with spectral
separation dependence is established~\cite{SDA}. The fractional
resolvent formula, rank-one inverse identity, determinant lemma,
Lanczos tridiagonalization and fast-subspace elimination used here
are classical. The repository's earlier path-transfer argument with
positive vertex masses is the starting point for
Lemma~\ref{lem:masspath}; it is repeated here to check the relative
metric and the exactly preserved leverage. The removal of arbitrary
vertex masses uses additional strongly weighted path blocks.

The present candidate contribution is the full determinant-admissible
relative spectral closure, its exact strict-count formula, and its
sharpness on paths with unit masses. The targeted source comparison
has not established historical priority. None of these statements is
an arbitrary quadratic-time matrix multiplication algorithm, a
finite-precision result, or a result for unweighted paths.

\begin{thebibliography}{9}
\bibitem{BCKS} B. Beckermann, A. Cortinovis, D. Kressner, M. Schweitzer.
Low-rank updates of matrix functions II: Rational Krylov methods.
SIAM J. Numer. Anal. 59 (2021), 1325--1347.
\url{https://arxiv.org/abs/2008.11501}.
\bibitem{SDA} S. Shmueli, P. Drineas, H. Avron.
Low-rank updates of matrix square roots.
Numer. Linear Algebra Appl. (2023), e2528.
\url{https://arxiv.org/abs/2201.13156}.
\bibitem{Repo} Fractional graph Laplacian updates without a spectral gap.
Repository working manuscript,
\texttt{matmul/graph-fractional-updates.tex}, exact path-transfer theorem.
\end{thebibliography}
\end{document}
